\section{Conclusion}\label{sec:conclusion}

We construct synthetic cultural agents by compiling the signs of declared aggregate preference anchors into group-indexed choice policies. A common bank of contrasting responses is labeled according to each group’s six-dimensional GPS sign profile, and a parameter-efficient adapter is fitted to those comparisons. Evaluating the resulting policies without country names in the prompt makes this population-specific labeling channel explicit and auditable. The fitted policy nevertheless remains jointly shaped by the pretrained model, the synthetic bank, and the complete sign profile.

The empirical application shows why anchor transfer and human agreement must be evaluated separately. The adapters learned the imposed synthetic comparisons, and their trust scores transferred the coarse GPS sign to new WVS wording. This pattern did not generalize uniformly across the remaining preference dimensions. Moreover, agreement with human country responses remained unresolved on the sixteen-country development panel. The central finding is therefore not that the adapters reproduce national cultures. It is that retaining a declared anchor and resembling human respondents are distinct empirical achievements: success on the former does not establish the latter.

The appropriate interpretation of these adapters is consequently as inspectable scoring instruments rather than synthetic populations. They provide a way to study whether a declared aggregate signal survives transport to new item text, while independent human evidence remains necessary to determine whether that signal corresponds to the intended construct. The next test is prospective: predeclare a set of candidate survey items, compare adapter-assisted and expert item selection, and evaluate both against a blinded human pilot. Such a study would determine whether anchored scoring improves survey pre-piloting rather than merely reproducing its inputs.
